\SourceLeaf{413}%
\par
I take as the axis of abscissas the straight line \(\mathrm{(II)}\), which intersects \((4),(5)\) at their midpoint \((6)\) at right angles. I count abscissas from \((6)\) as positive on the side of \((4),(5)\) where the angle from \((4)\) to \((5)\) appears less than \(180^0\), that is, on the right when the angles increase from left to right. Let ordinates be counted as positive in the direction from \((6)\) to \((5)\). On \(\mathrm{(II)}\) I label the points with abscissas
\[c.\operatorname{cotang}\vartheta=n-a,\quad c.\operatorname{cotang}\vartheta'=n,\quad c.\operatorname{cotang}\vartheta''=n+b\]
as \((1^*),(2^*),(3^*)\). These are the centers of the three circles that pass respectively through \((1),(2),(3)\) and all pass through \((4)\) and \((5)\). The radii of these circles are
\[\frac c{\sin\vartheta}=\BareRoot{(cc+(n-a)^2)},\quad\frac c{\sin\vartheta'}=\BareRoot{(cc+nn)},\quad\frac c{\sin\vartheta''}=\BareRoot{(cc+(n+b)^2)}\]
or, if we put \(\frac c{\sin\vartheta'}=r\), the other two become \(\BareRoot{(rr-2an+aa)},\BareRoot{(rr+2bn+bb)}\). Finally, let \((7)\) be the intersection of \(\mathrm{(I)}\) and \(\mathrm{(II)}\), \(T\) and \(t\) the distances of points \((2)\) and \((2^*)\) from \((7)\), and \(\varphi\) the angle between the corresponding rays of those lines, measured positively from \(\mathrm{(I)}\) to \(\mathrm{(II)}\) in the chosen sense. Thus the abscissa of \((7)=n-t\), so the coordinates of the three observation points are
\[\begin{array}{rcll}(1)&.\ .\ .\ .\ .&n-t+(T-A).\cos\varphi,&(T-A).\sin\varphi\\
(2)&.\ .\ .\ .\ .&n-t+T.\cos\varphi,&T.\sin\varphi\\
(3)&.\ .\ .\ .\ .&n-t+(T+B)\cos\varphi,&(T+B).\sin\varphi\end{array}\]
\par
The three unknown quantities \(t,T,\varphi\) are to be determined from the following equations, where for brevity \(x\) stands for \(\cos\varphi\).
\[\begin{aligned}tt+TT-2tTx&=rr\quad.\ .\ .\ .\ .\ .\ [1]\\
(t-a)^2+(T-A)^2-2(t-a)(T-A)x&=rr-2an+aa\\
(t+b)^2+(T+B)^2-2(t+b)(T+B)x&=rr+2bn+bb\end{aligned}\]
\par
In place of the last two equations I use the following ones, obtained by subtracting each from the first.
\[\begin{aligned}2at+2AT-2an-AA&=(2At+2aT-2aA)x\quad.\ .\ .\ [2]\\
2bt+2BT-2bn+BB&=(2Bt+2bT+2bB)x\quad.\ .\ .\ [3]\end{aligned}\]
